\documentclass[12pt,a4paper]{article}

% ---------- Paquetes ----------
\usepackage[utf8]{inputenc}
\usepackage[T1]{fontenc}
\usepackage[spanish]{babel}
\usepackage[margin=2.5cm, top=3cm, bottom=3cm]{geometry}
\usepackage{hyperref}
\usepackage{graphicx}
\usepackage{booktabs}
\usepackage{xcolor}
\usepackage{listings}
\usepackage{tikz}
\usepackage{float}
\usepackage{parskip}
\usepackage{fancyhdr}
\usepackage{array}
\usepackage{mdframed}
\usepackage{lmodern}
\usepackage[scaled=0.95]{helvet}
\usepackage{courier}
\renewcommand{\familydefault}{\sfdefault}
\usetikzlibrary{shapes.geometric, arrows.meta, positioning, fit, backgrounds}

% ---------- Colores de marca nexo ----------
\definecolor{nexoBlack}{HTML}{0A0E1A}
\definecolor{nexoCyan}{HTML}{00D4FF}
\definecolor{nexoMidnight}{HTML}{1A2744}
\definecolor{nexoGhost}{HTML}{F4F7FC}
\definecolor{nexoSlate}{HTML}{8A95A8}
\definecolor{nexoGreen}{HTML}{00E5A0}
\definecolor{nexoAmber}{HTML}{FFB800}
\definecolor{nexoRed}{HTML}{FF4D6A}
\definecolor{codeBg}{HTML}{F4F7FC}

% ---------- Listings ----------
\lstset{
  backgroundcolor=\color{codeBg},
  basicstyle=\ttfamily\footnotesize,
  keywordstyle=\color{nexoCyan!80!black}\bfseries,
  stringstyle=\color{nexoAmber!80!black},
  commentstyle=\color{nexoSlate}\itshape,
  breaklines=true,
  frame=single,
  framesep=5pt,
  rulecolor=\color{nexoSlate!30},
  showstringspaces=false,
  numbers=left,
  numberstyle=\tiny\color{nexoSlate},
  numbersep=8pt,
  tabsize=4,
  xleftmargin=1.5em,
}

\lstdefinelanguage{yaml}{
  keywords={true,false,null,yes,no},
  keywordstyle=\color{nexoCyan!80!black}\bfseries,
  sensitive=false,
  comment=[l]{\#},
  commentstyle=\color{nexoSlate}\itshape,
  stringstyle=\color{nexoAmber!80!black},
  morestring=[b]', morestring=[b]",
}

% ---------- Encabezado / pie de página ----------
\pagestyle{fancy}
\fancyhf{}
\renewcommand{\headrulewidth}{0pt}
\fancyhead[L]{%
  \begin{tikzpicture}[remember picture, overlay]
    \fill[nexoBlack] (current page.north west)
      rectangle ([yshift=-1.1cm]current page.north east);
    \node[anchor=west, text=white, font=\bfseries\small]
      at ([xshift=2.5cm, yshift=-0.55cm]current page.north west)
      {\textcolor{nexoCyan}{nexo} \textcolor{nexoSlate}{---}\ Infraestructura Backend};
    \node[anchor=east, text=nexoSlate, font=\small]
      at ([xshift=-2.5cm, yshift=-0.55cm]current page.north east)
      {Documento Interno · v3.0};
  \end{tikzpicture}%
}
\fancyfoot[C]{\small\textcolor{nexoSlate}{\thepage}}

% ---------- Hyperref ----------
\hypersetup{
  colorlinks=true,
  linkcolor=nexoCyan!60!black,
  urlcolor=nexoCyan!60!black,
  pdftitle={Infraestructura Backend — nexo courier},
  pdfauthor={nexo Engineering},
}

% ---------- Caja nexo ----------
\newmdenv[
  backgroundcolor=nexoBlack!5,
  linecolor=nexoCyan,
  linewidth=2pt,
  leftline=true, rightline=false, topline=false, bottomline=false,
  innerleftmargin=12pt, innerrightmargin=8pt,
  innertopmargin=6pt, innerbottommargin=6pt,
]{nexobox}

% ==========================================================
\begin{document}

% ---------- PORTADA ----------
\begin{titlepage}
  \pagecolor{nexoBlack}
  \color{white}
  \centering
  \vspace*{3cm}

  {\large\textcolor{nexoSlate}\itshape Tu puente entre dos mundos.\par}
  \vspace{1.2cm}

  {\fontsize{52}{56}\selectfont\bfseries
    \textcolor{white}{ne}\textcolor{nexoCyan}{x}\textcolor{white}{o}\par}
  \vspace{0.4cm}
  {\LARGE\bfseries Backend AWS --- Arquitectura Técnica\par}

  \vspace{1.8cm}
  \textcolor{nexoCyan}{\rule{0.55\textwidth}{1.5pt}}%
  \textcolor{nexoCyan}{\large\ $\bullet$}\par

  \vspace{1.5cm}
  {\large
    \textcolor{nexoSlate}{Empresa}\quad\textcolor{white}{nexo courier (Costa Rica)} \\[0.5em]
    \textcolor{nexoSlate}{Sistema}\quad\textcolor{white}{Autenticación, Base de Datos, Pedidos y API} \\[0.5em]
    \textcolor{nexoSlate}{Proveedor}\quad\textcolor{white}{Amazon Web Services (AWS)} \\[0.5em]
    \textcolor{nexoSlate}{Cuenta AWS}\quad\textcolor{white}{nexo-prod (769953010359)} \\[0.5em]
    \textcolor{nexoSlate}{Región}\quad\textcolor{white}{us-east-1} \\[0.5em]
    \textcolor{nexoSlate}{Versión}\quad\textcolor{white}{3.0} \\[0.5em]
    \textcolor{nexoSlate}{Fecha}\quad\textcolor{white}{Mayo 2026}
  }

  \vfill
  {\small\textcolor{nexoSlate}{Confidencial --- Uso Interno}}
\end{titlepage}

\pagecolor{white}
\color{black}

\tableofcontents
\newpage

% ==========================================================
\section{Resumen Ejecutivo}

\begin{nexobox}
El backend de \textbf{nexo courier} es 100\% serverless sobre AWS. No hay servidores
propios que mantener. Toda la infraestructura está definida como código con
\textbf{AWS CDK} y se despliega manualmente desde la máquina del equipo técnico.
\end{nexobox}

\vspace{0.5em}
El sistema maneja autenticación de usuarios mediante \textbf{AWS Cognito}, almacena
perfiles en \textbf{DynamoDB nexo-users}, pedidos en \textbf{DynamoDB nexo-orders},
direcciones de entrega en \textbf{DynamoDB nexo-addresses} y reseñas de clientes en
\textbf{DynamoDB nexo-reviews}, y expone un API protegido a través de
\textbf{5 Lambdas + API Gateway} con autorizador Cognito.
Todo dentro del \textit{Free Tier} de AWS para el volumen esperado en la fase inicial.

\subsection{Stack tecnológico}

\begin{table}[H]
\centering
\begin{tabular}{@{}lll@{}}
\toprule
\textbf{Componente} & \textbf{Tecnología} & \textbf{Proveedor} \\
\midrule
Autenticación & AWS Cognito User Pool & Amazon Web Services \\
Base de datos usuarios & AWS DynamoDB (on-demand) & Amazon Web Services \\
Base de datos pedidos & AWS DynamoDB (on-demand) & Amazon Web Services \\
Base de datos direcciones & AWS DynamoDB (on-demand) & Amazon Web Services \\
Lógica admin/usuarios & AWS Lambda (Node.js 20.x) & Amazon Web Services \\
Lógica pedidos & AWS Lambda (Node.js 20.x) & Amazon Web Services \\
Lógica direcciones & AWS Lambda (Node.js 20.x) & Amazon Web Services \\
Base de datos reseñas & AWS DynamoDB (on-demand) & Amazon Web Services \\
Lógica reseñas & AWS Lambda (Node.js 20.x) & Amazon Web Services \\
Lógica pagos (stub) & AWS Lambda (Node.js 20.x) & Amazon Web Services \\
Secretos Stripe & AWS Secrets Manager & Amazon Web Services \\
Servicio de email & AWS Lambda + Resend API & Amazon Web Services / Resend \\
Email transaccional & Resend (dominio nexocourier.com) & Resend \\
Trigger bienvenida & AWS Lambda (Cognito PostConfirmation) & Amazon Web Services \\
Email verificación & AWS Lambda (Cognito CustomMessage) + SES & Amazon Web Services \\
Secreto Resend & AWS Secrets Manager (nexo/resend) & Amazon Web Services \\
API & AWS API Gateway REST API + Cognito Authorizer & Amazon Web Services \\
IaC / Despliegue & AWS CDK (TypeScript) & Amazon Web Services \\
Alertas de costo & AWS Budgets & Amazon Web Services \\
\bottomrule
\end{tabular}
\caption{Stack tecnológico del backend}
\end{table}

% ==========================================================
\section{Arquitectura General}

\begin{lstlisting}[caption={Arquitectura general del backend}, numbers=none]
  Frontend (Next.js)
       |
       |-- signUp / confirmSignUp -------> Cognito User Pool
       |   (email verificacion branded)         |        |
       |                          CustomMessage |        | PostConfirmation
       |                                        v        v
       |                         nexo-cognito-custom-message
       |                         nexo-cognito-hooks -----> nexo-email-service
       |                                                         |
       |-- signIn ----------------------> Cognito User Pool      | Resend API
       |                                                         v
       |-- GET/POST /orders ----------> API Gateway --> Cognito Auth   email
       |-- PUT /admin/orders/{id}            |               |
       |   (cambia status)                   |          (valida JWT)
       |                                     v               v
       |            Lambda nexo-admin-users  Lambda nexo-orders-api (status change)
       |            Lambda nexo-addresses-api Lambda nexo-reviews-api      |
       |                     |                       |               fire-and-forget
       |            Cognito  DynamoDB      DynamoDB  |          nexo-email-service
       |           (admin)  nexo-users    nexo-orders|               |
       |                                    (+ totalBase,            v
       |                                     discountPct,       Resend API
       |                                     totalFinal)     email al cliente
\end{lstlisting}

\subsection*{Descripción de componentes}

\begin{description}
  \item[\textbf{Cognito User Pool}] Gestiona registro (con verificación de email),
    login y tokens JWT. El frontend se comunica directamente usando el SDK de Amplify v6.
  \item[\textbf{Cognito Authorizer}] Valida el JWT en cada request a API Gateway.
    Extrae los claims (sub, email, cognito:groups) y los pasa al Lambda.
  \item[\textbf{API Gateway}] Expone todos los endpoints REST. Todas las rutas
    requieren token JWT válido.
  \item[\textbf{Lambda nexo-admin-users}] CRUD de usuarios sobre Cognito y DynamoDB.
    Solo accesible por admins.
  \item[\textbf{Lambda nexo-orders-api}] Gestión de pedidos. Usuarios ven sus pedidos;
    admins ven todos, crean, actualizan estado, peso y total, y eliminan.
  \item[\textbf{Lambda nexo-addresses-api}] CRUD de direcciones de entrega CR. Máximo 2
    por usuario. Gestiona automáticamente la dirección predeterminada.
  \item[\textbf{DynamoDB nexo-users}] Perfiles extendidos de usuarios.
  \item[\textbf{DynamoDB nexo-orders}] Pedidos de envío registrados por los clientes.
  \item[\textbf{DynamoDB nexo-addresses}] Direcciones de entrega en Costa Rica.
\end{description}

% ==========================================================
\section{AWS Cognito User Pool}

\subsection{Identificadores}

\begin{table}[H]
\centering
\begin{tabular}{@{}ll@{}}
\toprule
\textbf{Campo} & \textbf{Valor} \\
\midrule
User Pool ID & \texttt{us-east-1\_Bphs6nkUf} \\
App Client ID & \texttt{5fqqogr9akgphaabk9jk4tl3ap} \\
Nombre & \texttt{nexo-users} \\
Región & \texttt{us-east-1} \\
ARN & \texttt{arn:aws:cognito-idp:us-east-1:769953010359:userpool/us-east-1\_Bphs6nkUf} \\
\bottomrule
\end{tabular}
\caption{Identificadores del User Pool}
\end{table}

\subsection{Configuración}

\begin{table}[H]
\centering
\begin{tabular}{@{}ll@{}}
\toprule
\textbf{Parámetro} & \textbf{Valor} \\
\midrule
Sign-in alias & Email (único por usuario) \\
Auto-verificación & Email (código de 6 dígitos al registrarse) \\
Recuperación de cuenta & Solo por email \\
Self sign-up & Habilitado \\
Removal policy & RETAIN \\
Remitente de emails & \texttt{notificaciones@nexocourier.com} (vía Amazon SES) \\
Nombre remitente & \texttt{nexo} \\
\bottomrule
\end{tabular}
\caption{Configuración general del User Pool}
\end{table}

\subsection{Triggers Lambda de Cognito}

\begin{table}[H]
\centering
\begin{tabular}{@{}lll@{}}
\toprule
\textbf{Trigger} & \textbf{Lambda} & \textbf{Descripción} \\
\midrule
PostConfirmation & \texttt{nexo-cognito-hooks} & Envía email de bienvenida al confirmar cuenta \\
CustomMessage & \texttt{nexo-cognito-custom-message} & Personaliza el email de verificación con template nexo \\
\bottomrule
\end{tabular}
\caption{Triggers Lambda del User Pool}
\end{table}

\begin{nexobox}
\textbf{PostConfirmation:} Dispara exactamente una vez, cuando el usuario confirma
su cuenta con el código de 6 dígitos. El Lambda \texttt{nexo-cognito-hooks} invoca
\texttt{nexo-email-service} con \texttt{type: 'WELCOME'} (await obligatorio ---
el contexto de Lambda se congela si no se espera).\\[4pt]
\textbf{CustomMessage:} Dispara antes de que Cognito envíe cualquier email del sistema
(verificación, reenvío de código). El Lambda construye un HTML con el template de nexo
y retorna el evento modificado. Cognito reemplaza \texttt{\{"\#\#\#\#"\}} con el código
real antes de enviar vía SES.
\end{nexobox}

\subsection{Amazon SES --- Envío desde Cognito}

Para que Cognito pueda enviar desde \texttt{notificaciones@nexocourier.com} (en lugar
del dominio genérico de Amazon), el dominio está verificado en \textbf{Amazon SES}
(\texttt{us-east-1}). Estado actual:

\begin{itemize}
  \item Dominio \texttt{nexocourier.com}: \textbf{Verified} + \textbf{DKIM Successful}
  \item Registros CNAME de DKIM agregados en Cloudflare (modo DNS only)
  \item Solicitud de acceso a producción (salir del sandbox) enviada a AWS
  \item En sandbox: solo envía a emails verificados en SES Console
\end{itemize}

\begin{nexobox}
\textbf{Importante:} SES y Resend son independientes. SES solo se usa para que
Cognito envíe el email de verificación desde el dominio propio. Resend se usa para
todos los emails transaccionales post-registro (bienvenida, status de pedidos).
\end{nexobox}

\subsection{Política de contraseñas}

\begin{table}[H]
\centering
\begin{tabular}{@{}ll@{}}
\toprule
\textbf{Requisito} & \textbf{Valor} \\
\midrule
Longitud mínima & 8 caracteres \\
Minúsculas & Requerido \\
Mayúsculas & No requerido \\
Números & Requerido \\
Símbolos & No requerido \\
\bottomrule
\end{tabular}
\caption{Política de contraseñas}
\end{table}

\subsection{Atributos del usuario}

\begin{table}[H]
\centering
\begin{tabular}{@{}llll@{}}
\toprule
\textbf{Atributo} & \textbf{Tipo} & \textbf{Requerido} & \textbf{Descripción} \\
\midrule
\texttt{email} & Estándar & Sí & Correo electrónico, mutable \\
\texttt{given\_name} & Estándar & Sí & Nombre, mutable \\
\texttt{family\_name} & Estándar & No & Apellido, mutable \\
\texttt{custom:tipo} & Personalizado & No & \texttt{persona} o \texttt{empresa} \\
\texttt{custom:movil} & Personalizado & No & Número de móvil \\
\texttt{custom:telefono} & Personalizado & No & Teléfono fijo \\
\bottomrule
\end{tabular}
\caption{Atributos del usuario en Cognito}
\end{table}

\begin{nexobox}
\textbf{Nota:} El atributo estándar \texttt{phone\_number} NO se usa. Los números
de contacto se almacenan en \texttt{custom:movil} y \texttt{custom:telefono} para
evitar el requisito de formato E.164 de Cognito.
\end{nexobox}

\subsection{Grupos}

\begin{table}[H]
\centering
\begin{tabular}{@{}lll@{}}
\toprule
\textbf{Grupo} & \textbf{Descripción} & \textbf{Permisos} \\
\midrule
\texttt{admin} & Administradores de nexo & CRUD usuarios + gestión de pedidos \\
\texttt{user} & Clientes registrados & Solo sus propios pedidos \\
\bottomrule
\end{tabular}
\caption{Grupos del User Pool}
\end{table}

\subsection{App Client}

\begin{description}
  \item[\textbf{Nombre}] \texttt{nexo-web}
  \item[\textbf{Client Secret}] No (SPA pública)
  \item[\textbf{Auth flows:}]
    \begin{itemize}
      \item \texttt{USER\_PASSWORD\_AUTH} --- login directo email/contraseña
      \item \texttt{USER\_SRP\_AUTH} --- Secure Remote Password
    \end{itemize}
\end{description}

% ==========================================================
\section{AWS DynamoDB}

\subsection{Tabla nexo-users}

\begin{table}[H]
\centering
\begin{tabular}{@{}ll@{}}
\toprule
\textbf{Campo} & \textbf{Valor} \\
\midrule
Nombre & \texttt{nexo-users} \\
Partition Key & \texttt{userId} (String) \\
Billing Mode & PAY\_PER\_REQUEST \\
Removal Policy & RETAIN \\
GSI & \texttt{email-index} (PK: email) \\
\bottomrule
\end{tabular}
\caption{Tabla nexo-users}
\end{table}

\begin{lstlisting}[language=bash, caption={Esquema de ítem en nexo-users}]
{
  "userId":    "string",    // sub de Cognito (PK)
  "email":     "string",
  "nombre":    "string",
  "apellido":  "string",
  "tipo":      "persona" | "empresa",
  "movil":     "string",
  "telefono":  "string",
  "createdAt": "string"    // ISO 8601
}
\end{lstlisting}

\subsection{Tabla nexo-orders}

\begin{table}[H]
\centering
\begin{tabular}{@{}ll@{}}
\toprule
\textbf{Campo} & \textbf{Valor} \\
\midrule
Nombre & \texttt{nexo-orders} \\
Partition Key & \texttt{orderId} (String, UUID) \\
Billing Mode & PAY\_PER\_REQUEST \\
Removal Policy & RETAIN \\
GSI & \texttt{userId-index} (PK: userId) \\
\bottomrule
\end{tabular}
\caption{Tabla nexo-orders}
\end{table}

\begin{lstlisting}[language=bash, caption={Esquema de ítem en nexo-orders}]
{
  "orderId":          "string",   // UUID (PK)
  "userId":           "string",   // sub de Cognito del dueno (GSI)
  "userName":         "string",   // para display en admin
  "userEmail":        "string",   // para notificaciones y display admin
  "trackingNumber":   "string",
  "description":      "string",
  "startDate":        "string",   // ISO 8601
  "status":           "en_ruta" | "atascado_aduana" | "bodega_cr"
                    | "pendiente_pago" | "pagado_en_ruta" | "entregado",
  "peso":             number,     // kg, lo carga el admin
  "totalPagado":      number,     // USD base (peso x tarifa), lo carga el admin
  "totalBase":        number,     // USD sin descuento — se calcula al llegar a bodega_cr
  "discountPct":      number,     // % Nexo Fiel aplicado (0, 3, 5 o 7)
  "totalFinal":       number,     // USD que el cliente debe pagar (con descuento)
  "updatedAt":        "string",   // ISO 8601
  // Campos de entrega — denormalizados al crear el pedido:
  "deliveryProvince": "string",   // opcional
  "deliveryCanton":   "string",   // opcional
  "deliveryDistrict": "string",   // opcional
  "deliverySenas":    "string"    // opcional
}
\end{lstlisting}

\begin{nexobox}
\textbf{Campos de pago (totalBase, discountPct, totalFinal):} Se calculan y persisten
automáticamente cuando el admin cambia el status a \texttt{bodega\_cr}. El Lambda
\texttt{nexo-orders-api} hace un Scan de todos los pedidos del usuario con status
billable (\texttt{bodega\_cr}, \texttt{pagado\_en\_ruta}, \texttt{entregado}), suma
los kg acumulados, aplica los tiers de Nexo Fiel y guarda los tres campos en DynamoDB.
Sirven como respaldo del cobro acordado y como input al servicio de pago futuro.
\end{nexobox}

\subsection{Tabla nexo-addresses}

\begin{table}[H]
\centering
\begin{tabular}{@{}ll@{}}
\toprule
\textbf{Campo} & \textbf{Valor} \\
\midrule
Nombre & \texttt{nexo-addresses} \\
Partition Key & \texttt{addressId} (String, UUID) \\
Billing Mode & PAY\_PER\_REQUEST \\
Removal Policy & RETAIN \\
GSI & \texttt{userId-index} (PK: userId) \\
\bottomrule
\end{tabular}
\caption{Tabla nexo-addresses}
\end{table}

\begin{lstlisting}[language=bash, caption={Esquema de ítem en nexo-addresses}]
{
  "addressId":  "string",    // UUID (PK)
  "userId":     "string",    // sub de Cognito del dueno (GSI)
  "province":   "string",    // provincia CR
  "canton":     "string",    // canton CR
  "district":   "string",    // distrito CR
  "senas":      "string",    // senas exactas de entrega
  "isDefault":  boolean,     // true si es la predeterminada
  "createdAt":  "string"     // ISO 8601
}
\end{lstlisting}

\begin{nexobox}
\textbf{Reglas de negocio:} Cada usuario puede tener máximo 2 direcciones. Siempre
debe existir exactamente una marcada como \texttt{isDefault = true}. Al eliminar la
dirección predeterminada, el Lambda asigna automáticamente \texttt{isDefault = true}
a la dirección restante. Al marcar una dirección como predeterminada, el Lambda
desactiva la otra (\texttt{isDefault = false}).
\end{nexobox}

\begin{nexobox}
\textbf{Separación de responsabilidades:} Las credenciales (contraseña hasheada,
tokens JWT) viven exclusivamente en Cognito. DynamoDB solo almacena datos de negocio.
Los pedidos de usuarios eliminados se \textbf{conservan} como registro histórico.
\end{nexobox}

\subsection{Tabla nexo-reviews}

\begin{table}[H]
\centering
\begin{tabular}{@{}ll@{}}
\toprule
\textbf{Campo} & \textbf{Valor} \\
\midrule
Nombre & \texttt{nexo-reviews} \\
Partition Key & \texttt{reviewId} (String, UUID) \\
Billing Mode & PAY\_PER\_REQUEST \\
Removal Policy & RETAIN \\
GSI & \texttt{userId-index} (PK: userId) \\
\bottomrule
\end{tabular}
\caption{Tabla nexo-reviews}
\end{table}

\begin{lstlisting}[language=bash, caption={Esquema de ítem en nexo-reviews}]
{
  "reviewId":  "string",   // UUID (PK)
  "userId":    "string",   // sub de Cognito del autor (GSI)
  "userName":  "string",   // nombre a mostrar publicamente
  "rating":    number,     // 1 a 5 (entero)
  "comment":   "string",   // minimo 10 caracteres
  "createdAt": "string"    // ISO 8601
}
\end{lstlisting}

\begin{nexobox}
\textbf{Reglas de negocio:} Cada usuario solo puede dejar \textbf{una} reseña.
El Lambda verifica vía GSI \texttt{userId-index} que no exista reseña previa del mismo
\texttt{userId} antes de insertar. Para poder dejar reseña el usuario debe tener al
menos un pedido con estado \texttt{entregado} en \texttt{nexo-orders}.
Los 3 ítems con \texttt{userId = "seed"} son datos de demostración insertados mediante
CDK Custom Resource al hacer deploy inicial (Ana Rodríguez, Carlos Jiménez, María González).
\end{nexobox}

% ==========================================================
\section{AWS Lambda}

\subsection{Lambda nexo-admin-users}

\begin{table}[H]
\centering
\begin{tabular}{@{}ll@{}}
\toprule
\textbf{Parámetro} & \textbf{Valor} \\
\midrule
Nombre & \texttt{nexo-admin-users} \\
Runtime & Node.js 20.x \\
Handler & \texttt{index.handler} \\
Timeout & 15 segundos \\
Variables de entorno & \texttt{USER\_POOL\_ID}, \texttt{TABLE\_NAME} \\
\bottomrule
\end{tabular}
\caption{Lambda de administración de usuarios}
\end{table}

\begin{table}[H]
\centering
\begin{tabular}{@{}lll@{}}
\toprule
\textbf{Método} & \textbf{Path} & \textbf{Descripción} \\
\midrule
GET & \texttt{/admin/users} & Lista todos los usuarios del User Pool \\
GET & \texttt{/admin/users/\{userId\}} & Obtiene un usuario por ID \\
PUT & \texttt{/admin/users/\{userId\}} & Actualiza atributos de un usuario \\
DELETE & \texttt{/admin/users/\{userId\}} & Elimina usuario de Cognito y DynamoDB \\
\bottomrule
\end{tabular}
\caption{Endpoints de nexo-admin-users}
\end{table}

\subsection{Lambda nexo-orders-api}

\begin{table}[H]
\centering
\begin{tabular}{@{}ll@{}}
\toprule
\textbf{Parámetro} & \textbf{Valor} \\
\midrule
Nombre & \texttt{nexo-orders-api} \\
Runtime & Node.js 20.x \\
Handler & \texttt{index.handler} \\
Timeout & 15 segundos \\
Variables de entorno & \texttt{TABLE\_NAME} (nexo-orders) \\
\bottomrule
\end{tabular}
\caption{Lambda de gestión de pedidos}
\end{table}

\begin{table}[H]
\centering
\begin{tabular}{@{}llll@{}}
\toprule
\textbf{Método} & \textbf{Path} & \textbf{Rol requerido} & \textbf{Descripción} \\
\midrule
POST   & \texttt{/orders} & Cualquier usuario auth & Crea pedido con tracking, descripción y dirección \\
GET    & \texttt{/orders} & Cualquier usuario auth & Lista pedidos propios (por userId del JWT) \\
GET    & \texttt{/admin/orders} & Grupo \texttt{admin} & Lista todos los pedidos (Scan) \\
POST   & \texttt{/admin/orders} & Grupo \texttt{admin} & Crea pedido para cualquier cliente \\
PUT    & \texttt{/admin/orders/\{orderId\}} & Grupo \texttt{admin} & Actualiza status, peso, totalPagado \\
DELETE & \texttt{/admin/orders/\{orderId\}} & Grupo \texttt{admin} & Elimina un pedido \\
\bottomrule
\end{tabular}
\caption{Endpoints de nexo-orders-api}
\end{table}

\subsection{Lambda nexo-addresses-api}

\begin{table}[H]
\centering
\begin{tabular}{@{}ll@{}}
\toprule
\textbf{Parámetro} & \textbf{Valor} \\
\midrule
Nombre & \texttt{nexo-addresses-api} \\
Runtime & Node.js 20.x \\
Handler & \texttt{index.handler} \\
Timeout & 15 segundos \\
Variables de entorno & \texttt{TABLE\_NAME} (nexo-addresses) \\
\bottomrule
\end{tabular}
\caption{Lambda de gestión de direcciones}
\end{table}

\begin{table}[H]
\centering
\begin{tabular}{@{}llll@{}}
\toprule
\textbf{Método} & \textbf{Path} & \textbf{Rol requerido} & \textbf{Descripción} \\
\midrule
GET    & \texttt{/addresses} & Cualquier usuario auth & Lista direcciones propias (por userId del JWT) \\
POST   & \texttt{/addresses} & Cualquier usuario auth & Crea dirección (máx. 2 por usuario) \\
PUT    & \texttt{/addresses/\{addressId\}} & Cualquier usuario auth & Edita dirección o cambia predeterminada \\
DELETE & \texttt{/addresses/\{addressId\}} & Cualquier usuario auth & Elimina dirección (solo si hay otra) \\
\bottomrule
\end{tabular}
\caption{Endpoints de nexo-addresses-api}
\end{table}

\begin{nexobox}
\textbf{Permisos IAM de nexo-addresses-api} sobre DynamoDB \texttt{nexo-addresses}:
acceso completo de lectura/escritura (PutItem, Query via GSI, UpdateItem, DeleteItem).
\end{nexobox}

\begin{nexobox}
\textbf{Verificación de admin en Lambda:} La función lee el claim
\texttt{cognito:groups} del JWT (disponible en
\texttt{event.requestContext.authorizer.claims}) y verifica que incluya
\texttt{'admin'} antes de ejecutar operaciones privilegiadas. Si no es admin,
retorna \texttt{403 Acceso denegado}.
\end{nexobox}

\subsection{Lambda nexo-reviews-api}

\begin{table}[H]
\centering
\begin{tabular}{@{}ll@{}}
\toprule
\textbf{Parámetro} & \textbf{Valor} \\
\midrule
Nombre & \texttt{nexo-reviews-api} \\
Runtime & Node.js 20.x \\
Handler & \texttt{index.handler} \\
Timeout & 15 segundos \\
Variables de entorno & \texttt{REVIEWS\_TABLE} (nexo-reviews), \texttt{ORDERS\_TABLE} (nexo-orders) \\
\bottomrule
\end{tabular}
\caption{Lambda de gestión de reseñas}
\end{table}

\begin{table}[H]
\centering
\begin{tabular}{@{}llll@{}}
\toprule
\textbf{Método} & \textbf{Path} & \textbf{Rol requerido} & \textbf{Descripción} \\
\midrule
GET    & \texttt{/reviews} & Público (sin auth) & Lista todas las reseñas, orden cronológico inverso \\
POST   & \texttt{/reviews} & Cualquier usuario auth & Crea reseña (con validaciones de negocio) \\
DELETE & \texttt{/admin/reviews/\{reviewId\}} & Grupo \texttt{admin} & Elimina una reseña \\
\bottomrule
\end{tabular}
\caption{Endpoints de nexo-reviews-api}
\end{table}

\begin{nexobox}
\textbf{Validaciones en POST /reviews:}
\begin{itemize}
  \item \texttt{rating} debe ser un entero entre 1 y 5.
  \item \texttt{comment} debe tener al menos 10 caracteres.
  \item El usuario debe tener al menos un pedido con estado \texttt{entregado}
    (consulta a \texttt{nexo-orders} vía GSI \texttt{userId-index}).
  \item Solo se permite una reseña por usuario (consulta a \texttt{nexo-reviews}
    vía GSI \texttt{userId-index}; retorna \texttt{409 Conflict} si ya existe).
\end{itemize}
\end{nexobox}

\begin{nexobox}
\textbf{Datos iniciales (seed):} Al hacer \texttt{cdk deploy} por primera vez se insertan
3 reseñas de demostración vía CDK Custom Resource (\texttt{AwsCustomResource}):
Ana Rodríguez (5\,★, nov.\ 2025), Carlos Jiménez (5\,★, dic.\ 2025) y María González
(5\,★, ene.\ 2026). Estos ítems usan \texttt{userId = "seed"} y no compiten con
usuarios reales.
\end{nexobox}

\subsection{Lambdas de Pagos (Stub)}

\begin{nexobox}
\textbf{Estado actual:} Las Lambdas de pago están desplegadas pero \textbf{no implementadas}.
Ambas retornan respuestas de marcador de posición. La integración real con Stripe está
pendiente de desarrollo.
\end{nexobox}

\begin{table}[H]
\centering
\begin{tabular}{@{}ll@{}}
\toprule
\textbf{Parámetro} & \textbf{Valor} \\
\midrule
Nombre & \texttt{nexo-payment-create} \\
Runtime & Node.js 20.x \\
Timeout & 30 segundos \\
Respuesta actual & \texttt{503 Service Unavailable} \\
Variables de entorno & \texttt{CORS\_ORIGINS} \\
\bottomrule
\end{tabular}
\caption{Lambda nexo-payment-create (stub)}
\end{table}

\begin{table}[H]
\centering
\begin{tabular}{@{}ll@{}}
\toprule
\textbf{Parámetro} & \textbf{Valor} \\
\midrule
Nombre & \texttt{nexo-payment-webhook} \\
Runtime & Node.js 20.x \\
Timeout & 30 segundos \\
Respuesta actual & \texttt{200 \{received: true\}} \\
Variables de entorno & \texttt{ORDERS\_TABLE} (nexo-orders), \texttt{SECRET\_ARN} \\
\bottomrule
\end{tabular}
\caption{Lambda nexo-payment-webhook (stub)}
\end{table}

\begin{table}[H]
\centering
\begin{tabular}{@{}llll@{}}
\toprule
\textbf{Método} & \textbf{Path} & \textbf{Rol requerido} & \textbf{Descripción} \\
\midrule
POST & \texttt{/payments/create}  & Cualquier usuario auth & Iniciar sesión de pago (stub: 503) \\
POST & \texttt{/payments/webhook} & Público (sin auth) & Recibir evento de Stripe (stub: 200) \\
\bottomrule
\end{tabular}
\caption{Endpoints de pagos (stub)}
\end{table}

\begin{nexobox}
\textbf{Integración Stripe planeada:}
\begin{itemize}
  \item Las credenciales se almacenan en AWS Secrets Manager bajo el nombre
    \texttt{nexo/stripe} con los campos \texttt{secretKey} y \texttt{webhookSecret}
    (actualmente con valor \texttt{PENDIENTE}).
  \item \texttt{nexo-payment-create} creará una Stripe Checkout Session y retornará
    la \texttt{checkoutUrl} al frontend.
  \item \texttt{nexo-payment-webhook} validará la firma del evento Stripe con
    \texttt{stripe.webhooks.constructEvent()} y actualizará el pedido al estado
    \texttt{pagado\_en\_ruta} en \texttt{nexo-orders}.
\end{itemize}
\end{nexobox}

\subsection{Lambda nexo-email-service}

\begin{table}[H]
\centering
\begin{tabular}{@{}ll@{}}
\toprule
\textbf{Parámetro} & \textbf{Valor} \\
\midrule
Nombre & \texttt{nexo-email-service} \\
Runtime & Node.js 20.x \\
Handler & \texttt{index.handler} \\
Timeout & 15 segundos \\
Variables de entorno & \texttt{RESEND\_SECRET\_ARN} \\
\bottomrule
\end{tabular}
\caption{Lambda centralizado de emails}
\end{table}

Lambda centralizado para todos los emails transaccionales. Cualquier otro Lambda
que necesite enviar un email lo invoca con \texttt{InvocationType: 'Event'}
(fire-and-forget). El Lambda lee la API Key de Resend desde Secrets Manager
(\texttt{nexo/resend}) con caché en memoria.

\begin{table}[H]
\centering
\begin{tabular}{@{}lll@{}}
\toprule
\textbf{type} & \textbf{Trigger} & \textbf{Descripción} \\
\midrule
\texttt{STATUS\_CHANGE} & Admin cambia status del pedido & Email de notificación con estado, tracking y card de pago \\
\texttt{WELCOME} & Cuenta confirmada (PostConfirmation) & Email de bienvenida con CTA a \texttt{/cuenta} \\
\bottomrule
\end{tabular}
\caption{Tipos de email del servicio}
\end{table}

\begin{nexobox}
\textbf{Email bodega\_cr:} Incluye una card de pago con fondo violeta (\texttt{\#F5F3FF}).
Si el usuario tiene descuento Nexo Fiel: muestra precio original tachado en gris,
badge ``Nexo Fiel X\%'' en morado y precio final en verde (\texttt{\#059669}).
Si no tiene descuento: solo el precio final sin tachado.\\[4pt]
\textbf{Statuses que generan email:} todos excepto \texttt{pendiente\_pago} (no usado
en el flujo actual). Los 6 estados configurados son: \texttt{en\_ruta},
\texttt{atascado\_aduana}, \texttt{bodega\_cr}, \texttt{pendiente\_pago},
\texttt{pagado\_en\_ruta}, \texttt{entregado}.
\end{nexobox}

\subsection{Lambda nexo-cognito-hooks}

\begin{table}[H]
\centering
\begin{tabular}{@{}ll@{}}
\toprule
\textbf{Parámetro} & \textbf{Valor} \\
\midrule
Nombre & \texttt{nexo-cognito-hooks} \\
Runtime & Node.js 20.x \\
Trigger & Cognito PostConfirmation \\
Timeout & 10 segundos \\
Variables de entorno & \texttt{EMAIL\_FUNCTION\_NAME} (\texttt{nexo-email-service}) \\
\bottomrule
\end{tabular}
\caption{Lambda de hooks de Cognito}
\end{table}

Dispara cuando el usuario confirma su cuenta (\texttt{PostConfirmation\_ConfirmSignUp}).
Invoca \texttt{nexo-email-service} con \texttt{type: 'WELCOME'} y devuelve el evento
original a Cognito (obligatorio). La invocación se hace con \texttt{await} para evitar
que el contexto de Lambda se congele antes de que el InvokeCommand se despache.

\subsection{Lambda nexo-cognito-custom-message}

\begin{table}[H]
\centering
\begin{tabular}{@{}ll@{}}
\toprule
\textbf{Parámetro} & \textbf{Valor} \\
\midrule
Nombre & \texttt{nexo-cognito-custom-message} \\
Runtime & Node.js 20.x \\
Trigger & Cognito CustomMessage \\
Timeout & 5 segundos \\
Variables de entorno & Ninguna \\
\bottomrule
\end{tabular}
\caption{Lambda de mensaje personalizado de Cognito}
\end{table}

Intercepta los triggers \texttt{CustomMessage\_SignUp} y \texttt{CustomMessage\_ResendCode}
para reemplazar el email de verificación genérico de Cognito con el template HTML
de nexo. Cognito sustituye \texttt{\{"\#\#\#\#"\}} con el código real antes de enviar
vía SES. Retorna siempre el evento original.

\subsection{Programa de Fidelidad --- Nexo Fiel}

Los descuentos se calculan automáticamente al cambiar status a \texttt{bodega\_cr},
sumando el peso de todos los pedidos del usuario en statuses billables:

\begin{table}[H]
\centering
\begin{tabular}{@{}lll@{}}
\toprule
\textbf{Kg acumulados} & \textbf{Descuento} & \textbf{Tier} \\
\midrule
$< 10$ kg & 0\% & Sin descuento \\
$\geq 10$ kg & 3\% & Nexo Fiel Bronce \\
$\geq 25$ kg & 5\% & Nexo Fiel Plata \\
$\geq 50$ kg & 7\% & Nexo Fiel Oro \\
\bottomrule
\end{tabular}
\caption{Tiers del programa Nexo Fiel}
\end{table}

Statuses que acumulan kg: \texttt{bodega\_cr}, \texttt{pagado\_en\_ruta}, \texttt{entregado}.
El pedido actual ya está en \texttt{bodega\_cr} cuando se hace el Scan, por lo que
sus kg están incluidos en el acumulado (correcto por diseño).

\subsection{Permisos IAM de las Lambdas}

\textbf{nexo-admin-users} sobre Cognito:
\begin{itemize}
  \item \texttt{cognito-idp:ListUsers}, \texttt{AdminGetUser}
  \item \texttt{cognito-idp:AdminUpdateUserAttributes}
  \item \texttt{cognito-idp:AdminDeleteUser}, \texttt{AdminDisableUser}
\end{itemize}

\textbf{nexo-admin-users} sobre DynamoDB \texttt{nexo-users}: acceso completo de lectura/escritura.

\textbf{nexo-orders-api} sobre DynamoDB \texttt{nexo-orders}: acceso completo de lectura/escritura (PutItem, Query, Scan, UpdateItem).

\textbf{nexo-reviews-api} sobre DynamoDB \texttt{nexo-reviews}: acceso completo de lectura/escritura (PutItem, Query vía GSI, Scan, DeleteItem).

\textbf{nexo-reviews-api} sobre DynamoDB \texttt{nexo-orders}: solo lectura (Query vía GSI \texttt{userId-index}, para verificar pedidos entregados).

\textbf{nexo-payment-create} sobre DynamoDB \texttt{nexo-orders}: solo lectura.

\textbf{nexo-payment-create} y \textbf{nexo-payment-webhook} sobre Secrets Manager \texttt{nexo/stripe}: solo lectura (\texttt{secretsmanager:GetSecretValue}).

\textbf{nexo-payment-webhook} sobre DynamoDB \texttt{nexo-orders}: acceso completo de lectura/escritura (para actualizar estado a \texttt{pagado\_en\_ruta}).

\textbf{nexo-email-service} sobre Secrets Manager \texttt{nexo/resend}: solo lectura (\texttt{GetSecretValue}).

\textbf{nexo-orders-api} sobre Lambda \texttt{nexo-email-service}: \texttt{lambda:InvokeFunction} (para invocación fire-and-forget al cambiar status).

\textbf{nexo-cognito-hooks} sobre Lambda \texttt{nexo-email-service}: \texttt{lambda:InvokeFunction} (para enviar email de bienvenida).

% ==========================================================
\section{AWS API Gateway}

\subsection{Configuración}

\begin{table}[H]
\centering
\begin{tabular}{@{}ll@{}}
\toprule
\textbf{Campo} & \textbf{Valor} \\
\midrule
Nombre & \texttt{nexo-admin-api} \\
Tipo & REST API \\
Stage & \texttt{prod} \\
URL Base & \texttt{https://qduhjqlmfd.execute-api.us-east-1.amazonaws.com/prod} \\
Autorizador & Cognito User Pools (\texttt{nexo-cognito-authorizer}) \\
\bottomrule
\end{tabular}
\caption{Configuración de API Gateway}
\end{table}

\subsection{Cognito Authorizer}

Todas las rutas requieren un token JWT válido en el header \texttt{Authorization}.
API Gateway valida la firma y expiración del token contra el User Pool antes de
invocar el Lambda. El Lambda recibe los claims directamente en
\texttt{event.requestContext.authorizer.claims}.

\subsection{CORS}

\begin{itemize}
  \item \textbf{Allow Origins:} \texttt{*}
  \item \textbf{Allow Methods:} GET, POST, PUT, DELETE, OPTIONS
  \item \textbf{Allow Headers:} Authorization, Content-Type
  \item \textbf{Preflight:} Manejado automáticamente (OPTIONS)
\end{itemize}

\subsection{Árbol de recursos}

\begin{lstlisting}[language=bash, caption={Estructura de recursos en API Gateway}, numbers=none]
/
├── orders/
│   ├── GET    --> nexo-orders-api  (pedidos del usuario autenticado)
│   └── POST   --> nexo-orders-api  (crear pedido con tracking y direccion)
│
├── addresses/
│   ├── GET    --> nexo-addresses-api  (direcciones del usuario)
│   ├── POST   --> nexo-addresses-api  (crear direccion, max 2)
│   └── {addressId}/
│       ├── PUT    --> nexo-addresses-api  (editar / cambiar predeterminada)
│       └── DELETE --> nexo-addresses-api  (eliminar, solo si hay otra)
│
├── reviews/
│   ├── GET    --> nexo-reviews-api  (todas las resenas, PUBLICO sin auth)
│   └── POST   --> nexo-reviews-api  (crear resena, auth requerido)
│
├── payments/
│   ├── create/
│   │   └── POST --> nexo-payment-create  (iniciar pago, auth requerido, stub 503)
│   └── webhook/
│       └── POST --> nexo-payment-webhook (evento Stripe, PUBLICO sin auth, stub 200)
│
└── admin/
    ├── users/
    │   ├── GET            --> nexo-admin-users  (listar usuarios)
    │   └── {userId}/
    │       ├── GET        --> nexo-admin-users  (obtener usuario)
    │       ├── PUT        --> nexo-admin-users  (actualizar usuario)
    │       └── DELETE     --> nexo-admin-users  (eliminar usuario)
    │
    ├── orders/
    │   ├── GET            --> nexo-orders-api   (todos los pedidos)
    │   ├── POST           --> nexo-orders-api   (crear pedido para cliente)
    │   └── {orderId}/
    │       ├── PUT        --> nexo-orders-api   (actualizar status/peso/total)
    │       └── DELETE     --> nexo-orders-api   (eliminar pedido)
    │
    └── reviews/
        └── {reviewId}/
            └── DELETE     --> nexo-reviews-api  (eliminar resena, solo admin)

NOTA: /reviews GET y /payments/webhook POST no requieren Authorization (publicos).
Todas las demas rutas requieren: Authorization: <JWT de Cognito>
\end{lstlisting}

% ==========================================================
\section{AWS CDK --- Infrastructure as Code}

\subsection{Stack}

\begin{table}[H]
\centering
\begin{tabular}{@{}ll@{}}
\toprule
\textbf{Campo} & \textbf{Valor} \\
\midrule
Stack Name & \texttt{NexoStack} \\
Cuenta & \texttt{769953010359} (nexo-prod) \\
Región & \texttt{us-east-1} \\
Stack ARN & \texttt{arn:aws:cloudformation:us-east-1:769953010359:stack/NexoStack/...} \\
\bottomrule
\end{tabular}
\caption{Datos del stack CDK}
\end{table}

\subsection{Outputs del stack}

\begin{table}[H]
\centering
\begin{tabular}{@{}ll@{}}
\toprule
\textbf{Output} & \textbf{Valor} \\
\midrule
\texttt{UserPoolId} & \texttt{us-east-1\_Bphs6nkUf} \\
\texttt{UserPoolClientId} & \texttt{5fqqogr9akgphaabk9jk4tl3ap} \\
\texttt{ApiUrl} & \texttt{https://qduhjqlmfd.execute-api.us-east-1.amazonaws.com/prod/} \\
\texttt{UsersTableName} & \texttt{nexo-users} \\
\texttt{OrdersTableName} & \texttt{nexo-orders} \\
\texttt{AddressesTableName} & \texttt{nexo-addresses} \\
\texttt{ReviewsTableName} & \texttt{nexo-reviews} \\
\bottomrule
\end{tabular}
\caption{Outputs del stack CDK}
\end{table}

\subsection{Comandos de despliegue}

\begin{lstlisting}[language=bash, caption={Despliegue manual de infraestructura}]
# Desde la carpeta infrastructure/
cd infrastructure/

# Ver cambios antes de aplicar (siempre recomendado)
npx cdk diff

# Desplegar
npx cdk deploy --require-approval never
\end{lstlisting}

\begin{nexobox}
\textbf{Política de deploy:} El despliegue de infraestructura es \textbf{manual y
deliberado}. No está integrado en el pipeline de CI/CD de Amplify. Siempre ejecutar
\texttt{cdk diff} antes de \texttt{cdk deploy} para revisar los cambios antes
de aplicarlos en producción.
\end{nexobox}

% ==========================================================
\section{AWS Budgets --- Alerta de Costos}

Se configura una alerta mensual de \textbf{\$20 USD} con dos notificaciones
al correo \texttt{nexxo.courier+prod@gmail.com}:

\begin{table}[H]
\centering
\begin{tabular}{@{}lll@{}}
\toprule
\textbf{Alerta} & \textbf{Umbral} & \textbf{Descripción} \\
\midrule
Aviso & 80\% (\$16 USD) & Gasto inusual, revisar servicios \\
Crítica & 100\% (\$20 USD) & Límite alcanzado \\
\bottomrule
\end{tabular}
\caption{Alertas de presupuesto mensual}
\end{table}

\begin{nexobox}
\textbf{Importante:} Los budgets de AWS \textbf{no bloquean} el gasto, solo notifican.
Para el volumen inicial de Nexo el costo real esperado es \$0--\$3/mes dentro del
Free Tier. La alerta en \$20 es una red de seguridad ante configuraciones incorrectas
o tráfico inesperado.
\end{nexobox}

% ==========================================================
\section{IAM y Seguridad}

\subsection{Usuario de despliegue}

\begin{table}[H]
\centering
\begin{tabular}{@{}ll@{}}
\toprule
\textbf{Campo} & \textbf{Valor} \\
\midrule
Usuario IAM & \texttt{nexobot-deploy} \\
Cuenta & \texttt{769953010359} (nexo-prod) \\
Uso & Deploy CDK desde máquina local \\
\bottomrule
\end{tabular}
\caption{Usuario IAM de despliegue}
\end{table}

\subsection{Configuración local de AWS CLI}

\begin{lstlisting}[language=bash, caption={Configurar credenciales}]
aws configure
# AWS Access Key ID:     [key de nexobot-deploy]
# AWS Secret Access Key: [secret de nexobot-deploy]
# Default region name:   us-east-1
# Default output format: json
\end{lstlisting}

% ==========================================================
\section{Estimación de Costos}

\begin{table}[H]
\centering
\begin{tabular}{@{}llll@{}}
\toprule
\textbf{Servicio} & \textbf{Free Tier} & \textbf{Costo excedente} & \textbf{Est. mensual} \\
\midrule
Cognito & 50,000 MAU & \$0.0055/MAU & \$0 \\
DynamoDB & 25 GB + 200M req & \$1.25/M writes & \$0 \\
Lambda & 1M req + 400K GB-s & \$0.20/millón req & \$0 \\
API Gateway & 1M requests/mes & \$3.50/millón & \$0 \\
AWS Budgets & 2 alertas gratis & \$0.10/alerta/mes & \$0 \\
\midrule
\multicolumn{3}{r}{\textbf{Total estimado fase inicial:}} & \textbf{\$0 USD} \\
\bottomrule
\end{tabular}
\caption{Estimación de costos mensual}
\end{table}

% ==========================================================
\section{Recuperación y Respaldo}

\begin{description}
  \item[\textbf{Cognito User Pool}] \texttt{removalPolicy: RETAIN} --- persiste si se
    destruye el CDK stack. Los usuarios siempre están seguros.
  \item[\textbf{DynamoDB nexo-users}] \texttt{removalPolicy: RETAIN} --- la tabla persiste
    independientemente del stack.
  \item[\textbf{DynamoDB nexo-orders}] \texttt{removalPolicy: RETAIN} --- los pedidos
    persisten aunque se elimine la cuenta del usuario.
  \item[\textbf{Código fuente}] Versionado en GitHub:
    \texttt{github.com/nexo-git/nexocode/infrastructure/}
  \item[\textbf{CloudFormation}] El estado del stack está en AWS y puede consultarse
    en cualquier momento desde la consola.
\end{description}

\begin{nexobox}
\textbf{Advertencia:} Nunca ejecutar \texttt{cdk destroy} en producción. A pesar del
\texttt{RETAIN} policy, los recursos auxiliares (roles IAM, API Gateway, Lambdas)
sí se eliminan y requieren re-deploy completo. Los datos en Cognito y DynamoDB
sobrevivirían, pero el sistema quedaría inoperativo.
\end{nexobox}

\vfill
\begin{center}
  \begin{tikzpicture}
    \fill[nexoBlack] (0,0) rectangle (\textwidth, 0.8cm);
    \node[anchor=west, text=nexoSlate, font=\small]
      at (0.3cm, 0.4cm) {nexo courier --- Confidencial --- Uso Interno --- v3.0};
    \node[anchor=east, text=nexoCyan, font=\small\itshape]
      at (\textwidth-0.3cm, 0.4cm) {Tu puente entre dos mundos.};
  \end{tikzpicture}
\end{center}

\end{document}
